\subsection{问题四：波动电价下重算两问}

\subsubsection{价格可观测性口径与预测器}

问题四的核心困难不在优化而在\textbf{信息}。附件 4 给出的是 2025 年逐日波动的实时电价，而“实时”意味着：当 0:00 制定当天计划时，当天 144 个时段的电价尚未发生。若把附件 4 直接当作已知输入，就等于假设电价可以被完全预知，这与题面不符，也会把问题四退化为一次“换电价重算”。因此本文明确采用如下口径：

\begin{quote}
\textbf{价格可观测性约定}：0:00 制定当天计划时，当天以及任何未来日的实时电价均不可观测；决策层只能使用历史电价（前一日及更早）与其他已有数据构造预测价；费用结算一律使用附件 4 的实际电价。
\end{quote}

在此约定下，链 \emph{4-3} 的调整层虽然拥有当天已实现时段的信息，6:00、12:00、18:00 时对当日剩余时段的电价\textbf{仍然不可知}，只能用滚动更新的预测；只有结算时才使用实际价。主预测器取“前一日同时段”$\hat p^{(0)}_{d,i}=p^{act}_{d-1,i}$：该预测器不需要估计参数，且与附件 4 的日内形状稳定性相容。为保证结论不依赖预测器选择，本文同时给出“水平 $\times$ 时段”双因子预测器（用最近 7 个可用历史日的日均价作为当日水平、用历史同时段的归一化形状作为日内结构）与“历史同时段均值”预测器作为基准对照。链 \emph{4-3} 的调整层在此基础上做滚动水平校正：

\begin{equation}
\hat p^{(m)}_{d,i}=
\begin{cases}
p^{act}_{d,i}, & i\le 6m\ \text{（已实现，可直接观测）}\\
\kappa_m\,\hat p^{(0)}_{d,i}, & i>6m
\end{cases},
\qquad
\kappa_m=\mathrm{clip}\!\left(A_m/B_m,\,0.5,\,2.0\right),\quad \kappa_0\equiv1
\end{equation}

其中 $A_m$ 为当天截至时刻 $m$ 已实现电价的均值，$B_m$ 为前一日同时段均价；含义是：若当天价格整体高于前一日同日同时段水平，就把剩余时段的预测按比例上调。实测 $\kappa_m$ 的取值区间为 $[0.66233,1.54939]$，从未触及上下界，即校正始终是内点解。实时电价预测与“用预测价做决策、用实际价结算”的分工在电力市场与储值优化文献中已有共识\cite{ref-lago-2018-spot-price-deep-learning,ref-myttenaere-2016-mape-regression,ref-chitsaz-2018-price-forecast-btm-storage,ref-khani-2015-online-adaptive-rtod,ref-staffell-2016-maximising-storage-value,ref-mathaba-2014-price-forecast-benefit,ref-hodge-2018-forecast-value-storage}。

预测器的滚动回测结果（只用被测日之前的信息、覆盖交付期 334 天）列于表~\ref{tab:p4fc}。

\begin{table}[htbp]
\centering
\small
\caption{价格预测器的滚动回测（交付期 334 天，单位：元/kWh 或百分数）}
\label{tab:p4fc}
\begin{tabular}{lrrrrrr}
\toprule
预测器 & MAE & MAPE & RMSE & $P_{50}$ & $P_{90}$ & $P_{99}$ \\
\midrule
前一日同时段（主口径） & 0.084326 & 13.7801\% & 0.123837 & 0.0500 & 0.2294 & 0.3773 \\
水平 $\times$ 时段双因子 & 0.093387 & 15.2904\% & 0.121566 & 0.0723 & 0.206537 & 0.335853 \\
历史同时段均值 & 0.100231 & 16.3301\% & 0.136878 & 0.0710 & 0.241094 & 0.396309 \\
\bottomrule
\end{tabular}
\end{table}

表~\ref{tab:p4fc} 提示一个必须成对报告的现象：在三个基准预测器中，主口径的\textbf{平均类误差}最小（MAE 与 MAPE 均为三者最优），而双因子预测器的\textbf{尾部类误差}最小（RMSE、$P_{90}$、$P_{99}$ 均优于其余两个）。两个预测器的优劣方向相反，因此只报单项指标会得出误导性结论，本文按平均类与尾部类指标成对给出；作为消融对照的自回归预测器在 MAE（0.080664）上低于主口径，本文如实披露该反例，但不据此替换主口径。另外，这两个基准预测器在决策费用上给出的结果几乎相同，说明“预测机制”在本题中只提供一条独立的信息轴，不应被计为两个独立的敏感性来源。

\begin{figure}[htbp]
\centering
\includegraphics[width=0.80\textwidth]{fig15_p4_backtest.png}
\caption{价格预测器的滚动回测：平均类与尾部类指标的最优方向相反}
\label{fig:p4fc}
\end{figure}

\subsubsection{两条链的模型}

\textbf{链 \emph{4-2}（重算问题二）}。由于 0:00 不知道当天电价，链 \emph{4-2} 不能用“单一全年长时域线性规划”——那需要整年的价格序列。本文采用逐日向前递推：对每个交付日 $d$，用预测价与当天已知的实际负载、实际光伏求解一个单日规划

\begin{equation}
\min\ \sum_{i}\hat p^{(0)}_{d,i}b_{d,i}+\sum_{i}\alpha_{em}\hat p^{(0)}_{d,i}q^{em}_{d,i}
\quad\text{s.t.}\quad
b_i+q^{em}_i+PV^{act}_i\Delta t+q^{dis}_i=L_i\Delta t+c_i+s_i
\end{equation}

提交计划量后断言 $q\equiv b$，储电量按日边界连续性传递、终端自由。全年 365 次单日规划共约 11.9 s。

\textbf{链 \emph{4-3}（重算问题三）}。在问题三的三层结构上，把计划层的价格换成 $\hat p^{(0)}$（$\kappa_0\equiv1$）、调整层的价格换成式 (16) 的 $\hat p^{(m)}$，其余（预报降尺度、偏差分段、结算层闭式）与问题三相同。全年共 1460 个单日层，约 38.9 s。

\textbf{决策—结算分离与“价格预报误差的代价”}。同一组提交的决策量可以用两套价格分别计价：按决策层信念价计得 $C^{fc}$，按附件 4 实际价计得 $C^{act}$。本文以 $C^{act}$ 为唯一主结果，并把差额 $\Delta C^{price}=C^{act}-C^{fc}$ 定义为“价格预报误差的代价”单独报告。这一分解与问题三中“预报误差的代价”在结构上同构：一个是价格预测错了的代价，一个是光伏预测错了的代价。日前与实时市场的两结算与价格不确定性对储能收益的影响在文献中已被系统研究\cite{ref-krishnamurthy-2018-arbitrage-da-rt,ref-yang-2023-da-rt-shared-storage,ref-bottieau-2020-imbalance-settlement-forecast,ref-yang-2020-deviation-penalty-distribution,ref-paulauskas-2024-battery-scheduling-arbitrage}。

\subsubsection{主结果与三分项归因}

两条链的交付期主结果列于表~\ref{tab:p4main}。链 \emph{4-2} 的交付期费用为 $13\,006\,411.041148$ 元，链 \emph{4-3} 为 $15\,289\,050.714738$ 元，后者比前者高

\begin{equation}
2\,282\,639.673590=100\,230.446225+522\,512.788776+1\,659\,896.438588\quad(\text{单位：元})
\end{equation}

即差额可以\textbf{完全分解}为三项：计划口径差（两链计划层的价格与信息集不同导致的计划费用差异）、偏差结算制度（链 \emph{4-3} 独有的 $0.5\times/1.5\times$ 双向分段费用）与紧急购电（光伏预报误差在结算层形成的缺口），恒等式残差不超过 $10^{-6}$ 元（来自表格的六位小数舍入）。这一“三分项归因”是本文对问题四最主要的定量结论，它把一个 $2.28\times10^{6}$ 元的总差额变成了三条各自可解释、可单独核验的通道。

\begin{table}[htbp]
\centering
\small
\caption{问题四两条链的交付期主结果（334 天；费用单位：元，电量单位：kWh）}
\label{tab:p4main}
\begin{tabular}{p{5.4cm}rr}
\toprule
量 & 链 \emph{4-2} & 链 \emph{4-3} \\
\midrule
计划购电费 $C_{\text{plan}}^{act}$ & 13\,006\,411.041148 & 13\,106\,641.487373 \\
调整结算费 $C_{\text{adj}}^{act}$ & 0（无该项） & 522\,512.788776 \\
紧急购电费 $C_{\text{em}}^{act}$ & 0（机制性恒零） & 1\,659\,896.438588 \\
\textbf{总费用 $C_{\text{total}}^{act}$} & \textbf{13\,006\,411.041148} & \textbf{15\,289\,050.714738} \\
全期（365 天）总费用 & 14\,755\,884.322036 & 17\,181\,202.515506 \\
信念价口径费用 $C^{fc}$ & 12\,485\,620.202113 & 15\,030\,565.249034 \\
价格预报误差的代价 $\Delta C^{price}$ & 520\,790.839035 & 258\,485.465704 \\
计划购电量 $\sum b$ / 最终购电量 $\sum q$ & 20\,241\,516.374467 / 20\,241\,516.374467 & 20\,411\,764.664671 / 20\,807\,826.888962 \\
紧急购电量 $\sum q^{em}$ / 弃光量 $\sum s'$ & 0 / 990\,168.090312 & 394\,026.040750 / 1\,858\,369.774422 \\
净负荷（两链逐位相同） & 17\,939\,189.494967 & 17\,939\,189.494967 \\
\bottomrule
\end{tabular}
\end{table}

量级与上界自检：两链交付期费用均落在 $[3\times10^{6},5\times10^{7}]$ 元带内（日均分别为 38941.35 元与 45775.60 元）；用“完全不用储能”的可行策略构造的价格项上界为 17\,026\,782.858099 元与 18\,884\,193.220703 元，两链主值均低于上界，对应的储能价值为 4\,020\,371.82 元与 1\,737\,732.14 元。注意不能用附件 4 的最小电价构造下界——由于存在 0.0076 元/kWh 的极端低值，该下界不足实际费用的 1.2\%，不具判别力。

表~\ref{tab:p4a} 给出两条链在题面四个指定日期的结果。链 \emph{4-2} 报的是计划购电量 $b$、全天购电费等于计划费；链 \emph{4-3} 报的是最终购电量 $q$、全天购电费为三项之和。两链的数值必然不同（价格情景与决策结构均不同）：四个日期的紧急购电在链 \emph{4-2} 上全为零、在链 \emph{4-3} 上分别为 5、7、9、4 个合并区间。以 2025 年 3 月 20 日为例，链 \emph{4-2} 的六个 4 小时块充电量为 $[8082.85,3333.33,5658.77,5322.37,0.00,833.33]$ kWh、放电量为 $[0,6469.73,2777.38,254.72,6478.06,2836.94]$ kWh；链 \emph{4-3} 的充电量为 $[8082.85,3333.33,2856.65,6959.71,522.57,833.33]$ kWh、放电量为 $[0,6417.07,2313.92,250.65,6478.06,2836.94]$ kWh。两链该日的端部储电量均为 1200.0 kWh，属逐日滚动切片，与问题一的单日周期端点（恒 6000 kWh）语义不同。

\begin{table}[htbp]
\centering
\footnotesize
\caption{问题四：四个指定日期的指定时段购电量与全天合计（购电量单位：kWh，购电费单位：元）}
\label{tab:p4a}
\begin{tabular}{llrrrrrrrr}
\toprule
链 & 日期 & 10:00 & 12:00 & 14:00 & 16:00 & 18:00 & 20:00 & 全天购电量 & 全天购电费 \\
\midrule
\emph{4-2} & 2025-03-20 & 0.00 & 558.311667 & 0.00 & 507.324167 & 0.00 & 726.368333 & 66\,622.786007 & 41\,543.092634 \\
 & 2025-06-21 & 0.00 & 0.00 & 0.00 & 27.005000 & 0.00 & 0.00 & 32\,532.468316 & 15\,588.839124 \\
 & 2025-09-23 & 0.00 & 380.021250 & 0.00 & 493.582500 & 2.645000 & 61.911667 & 66\,265.421275 & 43\,331.557914 \\
 & 2025-12-21 & 0.00 & 906.683333 & 0.00 & 795.930417 & 712.633333 & 30.268333 & 94\,061.221951 & 70\,297.812713 \\
\emph{4-3} & 2025-03-20 & 0.00 & 553.563367 & 0.00 & 641.129717 & 0.00 & 726.368333 & 67\,656.452692 & 52\,443.510373 \\
 & 2025-06-21 & 0.00 & 0.00 & 0.00 & 112.922733 & 0.00 & 0.00 & 34\,007.160791 & 19\,429.156761 \\
 & 2025-09-23 & 0.00 & 358.668517 & 0.00 & 550.185183 & 2.645000 & 61.911667 & 68\,433.174459 & 50\,189.088074 \\
 & 2025-12-21 & 0.00 & 833.333333 & 0.00 & 806.735183 & 712.633333 & 30.268333 & 94\,508.735661 & 76\,068.603362 \\
\bottomrule
\end{tabular}
\end{table}

\subsubsection{机制解释}

\textbf{链 \emph{4-2} 的价格不确定性不产生能量缺口}（图~\ref{fig:p4a}）。决策层使用的信念价与结算层的实际价之间存在系统偏差，但链 \emph{4-2} 中负载与光伏在 0:00 即完全已知，价格预测错了只会让计划买贵或买便宜、不会造成供给不足，因此紧急购电恒为零（实测最大绝对值 $2.27\times10^{-13}$ kWh，属求解器容差级噪声，故表述为“机制性恒零”而非“逐位为零”）。价格误差的全部后果体现为 520\,790.839035 元的费用差。

\textbf{链 \emph{4-3} 的缺口来自光伏预报误差}。该链的紧急购电全部由“计划层与调整层用附件 3 预报、结算层用附件 2 实际”这一信息落差造成，与价格预测器无关：把 $\kappa_m$ 冻结为 1（即不做滚动校正）对实际费用几乎没有影响（相差 $1.49\times10^{-8}$ 元），但会显著改变信念价口径下的费用差（由 258\,485.47 元升至 573\,883.11 元）。这说明滚动价格校正的作用在于改善“决策者对账单的预期”，而非降低实际支出。

\textbf{两链的相对敏感性方向一致但幅度不同}。在 10 个共同情景上两链的费用变化方向同号；但在价格、负载与预测器三类扰动上，链 \emph{4-3} 反而比链 \emph{4-2} \emph{更不敏感}（分别低约 10\%、17\%、24\%），机制是调整层用当天已实现电价做了水平校正，形成部分对冲；相应地，链 \emph{4-3} 对光伏相关扰动远为敏感（为链 \emph{4-2} 的 42--65 倍）。这一“一边更钝、一边更锐”的对比正是三条归因通道的直接体现。

\begin{figure}[htbp]
\centering
\includegraphics[width=0.80\textwidth]{fig14_p4_price.png}
\caption{决策层信念价与结算层实际价的分离（两链共用）}
\label{fig:p4a}
\end{figure}

\begin{figure}[htbp]
\centering
\includegraphics[width=0.80\textwidth]{fig16_p4_cost.png}
\caption{两链费用分项与交付期差额的三分项归因}
\label{fig:p4b}
\end{figure}

\subsubsection{稳健性与消融}

\textbf{鲁棒性}。链 \emph{4-2} 的 162 个情景全部成功、全部判据通过，分级为“稳定”；其最大敏感源是题面硬参数而非价格：往返效率的 $-20\%$ 档使费用上升 13.345\%，储电量上限收紧 20\% 使费用上升 5.141\%，而紧急购电倍数在六档上给出\textbf{完全相同}的费用（因紧急购电恒为零，该参数在本链上是空参数），初始储电量对交付期费用没有响应。链 \emph{4-3} 的 198 个情景中 195 个成功，3 个不可行情景全部是预注册的购电上限负对照（3500、4000、4500 kW）；另有 1 条判据未通过：光伏与负载的白噪声 $\sigma=10\%$ 族使费用均值上升 10.5919\%，超过 5\% 的预注册阈值，机制是紧急购电项的 $\max(0,\cdot)$ 凸性系统性抬高期望紧急购电，而所有受判样本均可行、恒等式残差为零，属结构性而非实现缺陷。据此链 \emph{4-3} 分级为“条件稳定（需给出适用边界）”：在 $\sigma\le5\%$ 的输入扰动下结论稳健，$\sigma=10\%$ 属超设计工况；结论\textbf{对价格侧稳健、对光伏侧不稳健}。本文不通过调参或放宽判据使该判据“通过”。

\textbf{消融与实现正确性锚点}。三组决定性对照（13 个对照项分链展开为 23 个算例）全部执行完毕。

第一，\textbf{退化情景复现（实现正确性锚点）}：把链 \emph{4-3} 的价格数据整体替换为附件 1 的单日曲线（逐日重复 365 次），其余输入与解的选择规则不变，此时“前一日同时段”预测在全年上精确成立、滚动校正因子恒为 1，模型理应退化为问题三。实测六项指标\textbf{逐位复现}问题三结果（全期 16\,179\,176.145228 元、交付期 14\,540\,616.335652 元、计划费 13\,885\,669.475555 元、调整费 563\,098.132533 元、紧急购电费 1\,730\,408.537141 元、交付期紧急购电量 391\,274.485914 kWh），六项相对差均不超过 $4.86\times10^{-13}$，远小于 $10^{-9}$ 的判据。这是本文对“问题四的实现是否正确”给出的最硬证据：一条完全独立于价格情景的退化路径必须精确回到已知答案。

第二，\textbf{两阶段随机规划对照（是否值得引入随机优化）}。为回答“点预测加确定性优化是否足够”，本文构造两阶段随机线性规划：第一阶段选择全天共用的计划购电量，第二阶段在每个价格情景内调整并结算；情景取自附件 4 的历史实际价格（对目标日取此前同型日的 144 点实际价向量、等概率，共 24 个情景），并裁剪为 12 个代表日的确定性等价大规划。结果显示，链 \emph{4-2} 上随机方案使同一窗口内的期望费用比点预测方案低 3.2860\%（13\,412.83 元），计划量的形状差异表现为“0--4 时与 16--20 时增购、4--8 时与 20--24 时减购”，即在低价时段多买、高价时段少买，属数量对冲。链 \emph{4-3} 上随机方案报告的降幅为 19.02\%，但这是因为预注册的追索结构使偏差变量恒为零、其期望费用与链 \emph{4-2} 逐位相同，因此\textbf{不得}把该降幅表述为“对冲在链 \emph{4-3} 上有 19\% 的收益”。综合判断：随机优化确实有收益，但在 12 个代表日窗口上不超过 3.3\%，且需要额外的场景建模与计算预算，本文据此维持“价格点预测 + 确定性线性规划”为主口径；这一判断也与“随机或鲁棒优化必须由问题结构与预算证明必要性”的方法论要求一致\cite{ref-akbari-2019-stochastic-robust-arbitrage,ref-oconnor-2025-conformal-price-forecast,ref-rawa-2023-mg-stochastic}。

第三，\textbf{滚动时域对照（前瞻深度）}。在决策时刻与承诺规则完全不变的前提下，只改变“优化时看多远”，把计划层的前瞻深度设为 $H\in\{1,3,7,14\}$ 天，窗口内的未来变量只用于优化、不进入承诺。$H=1$ 必须逐位复现主口径，实测全期与交付期费用差均为 0、整条储电量轨迹逐点差为 0，构成基线闸门。$H\ge3$ 时两链交付期费用分别再降 0.1785\% 与 0.1492\%，且前瞻深度在 $H=3$ 之后饱和（$H$ 由 3 增到 14 的费用增量仅 $-3.58$ 元与 $-5.49$ 元）。更重要的是，该对照直接回答了前文提出的结构性问题：主口径与 $H=1$ 下两链的日末储电量全部贴住下限 1200 kWh（364/365 与 365/365 天），而 $H\ge3$ 时日末储电量均值升至约 5335.9 与 5327.9 kWh、贴下限日数降至 30/365、唯一取值数增至 190 与 191。因此\textbf{“日边界不携带电量”只是“逐日 0:00 决策加终端自由”这一结构下的性质，不是与模型无关的普遍规律}，本文不作外推。此外如实登记一个反例：把两阶段随机与前瞻深度组合（24 个场景 $\times$ $H=3$、12 个代表日）后，期望费用反而比同窗点预测方案高 2.6959\%（11\,004.04 元），即“对冲加前瞻”在该窗口上\textbf{不叠加}收益；该结论只对窗口成立，本文不据此修改预注册设置。

最后，本文明确维持两条链的主口径不变：链 \emph{4-2} 为“前一日同时段预测价 + 逐日确定性线性规划”，链 \emph{4-3} 为“前一日同时段预测价 + 滚动水平校正 + 三层递推 + 双向分段结算”。模型族对照对费用结论的改变不超过 0.28\% 与 0.15\% 且无方向反转，退化情景锚点提供了实现正确性的硬证据，随机对照的收益幅度不足以证明更换主口径的必要性。
